In a scanning tunnelling microscope, a metal tip is brought within about \SI{1}{\nano\m} of a surface; a small voltage makes electrons tunnel across the gap. For the electrons, the gap is a barrier about as high as the work function, $V_0-E \approx \dEO$. The current is proportional to the fraction of electrons that tunnel, $T \approx e^{-2\kappa d}$, where $d$ is the width of the gap.

\begin{abcliste}
\abc Calculate the decay constant $\kappa$.

\abc By what factor does the current grow when the tip comes $\Delta d = \ddO$ closer?

\abc Why does the microscope show single atoms?
\end{abcliste}
